\section{Related Work}

\noindent \textbf{Latency in LLM-based Agent} ~ LLM-based agents have gained widespread adoption in real-world applications \citep{ge2024openagi, wu2023autogen} due to their strong reasoning and planning capabilities. However, significant latency issues have been documented across multiple domains: \cite{jaech2024openai} reports allocating 24-hour processing windows for o1-based agents to complete machine learning engineering tasks, while \cite{elfleet2024investigating} examines how perceived agent latency impacts user experience in virtual reality environments. \cite{liu2023llm} further notes that the extended inference times of LLM-based agents render them unsuitable for highly interactive and real-time applications such as gaming.

Several research directions have emerged to address these latency challenges. System 1.x \citep{saha2025system} leverages dual-process theory of cognition \citep{evans2003two, hua2022system}, dynamically substituting the primary backbone with smaller models when planning subtasks are identified as straightforward, thereby reducing individual inference times. EcoAct \citep{zhang2024ecoact} addresses a different source of delay by optimizing prompt length through the selective omission of tool descriptions irrelevant to the current step. Interactive Speculative Planning (ISP) \citep{hua2024interactive} accelerates execution by parallelizing action generation and verification, transforming the inherently autoregressive planning pipeline to non-autoregressive. However, a critical gap exists in the literature: these approaches either compromise performance guarantees to achieve acceleration, require extensive pre-deployment preparation, or incur prohibitive computational costs without providing mechanisms to control latency trade-offs. Our work bridges this gap by \textbf{introducing a lossless self-optimizing, zero-setup framework} based on speculative execution that can adaptively balance latency reduction against operational costs, providing the first comprehensive solution for deploying accelerated LLM agents in diverse real-world contexts with varying efficiency requirements.

\noindent \textbf{Speculative execution} ~ \citep{kocher2020spectre, xiao2019speechminer} originated in computer architecture as a technique to reduce latency by executing instructions before their necessity is confirmed, with mechanisms to roll back incorrect speculations. %Classical applications include branch and value prediction \citep{gabbay1996speculative} in superscalar processors, aggressive memory prefetching \citep{pugsley2014sandbox}, and optimistic concurrency control in database systems—all leveraging otherwise idle computational resources to increase effective parallelism while preserving correctness guarantees.
This principle has been successfully adapted to decoder-only LLMs for autoregressive token generation. Speculative decoding \citep{zhang2024beyond, leviathan2023fast} accelerates text generation through a two-model approach: an efficient draft model predicts multiple future tokens, which a more capable while slow target model verifies in parallel. Tokens that pass verification are committed to the output, while rejected tokens are regenerated. Research consistently demonstrates that the latency of speculative pipelines depends critically on the alignment between draft and target components \citep{bachmann2025judge}. Consequently, both offline knowledge-distillation schemes and online learning methods have been developed
%, thereby increasing token acceptance rates and further reducing latency. The literature
including numerous refinements such as adaptive draft lengths \citep{liupearl}, ensemble-draft strategies \citep{lee2024inbatch, fu2025speculative}, and online distillation methods \citep{liu2023online, zhou2023distillspec, ouyang2024temperature, gui2024boosting} that continuously align draft models with target models, demonstrating that speculation can reduce decoding latency while maintaining lossless generation distribution. ISP extends this paradigm to agent planning, where an approximation agent $(\approxi)$ generates tentative planning steps while a stronger target agent $(\target)$ verifies them asynchronously. While ISP's plug-and-play approach preserves planning performance, its fixed speculation step leads to either excessive costs or insufficient acceleration. Our work addresses this limitation by developing a dynamic speculation system that automatically determines optimal steps through online reinforcement learning, significantly reducing costs while preserving speed. Additionally, we introduce user-controllable mechanisms that expose the latency-cost continuum, enabling precise system calibration based on operational priorities, a capability absent in all previous speculative execution approaches.


%Consequently, an unnecessary approximation agent step can incur costs orders of magnitude higher than a misplaced token-level verification. 
% Our work addresses speculation step control in the agent-planning context by developing a cost-aware controller that maintains the latency advantages of speculative planning while substantially reducing the monetary overhead unique to multi-step action generation. Crucially, our controller enables user-specified preferences along the cost-time saving continuum, allowing dynamic adaptation to varying operational requirements and economic constraints.
